\documentclass[11pt]{article}
\usepackage[a4paper,margin=25mm]{geometry}
\usepackage{amsmath,amssymb,mathtools}
\usepackage{fontspec,unicode-math}
\setmainfont{Latin Modern Roman}
\setsansfont{Latin Modern Sans}
\setmonofont{Latin Modern Mono}
\setmathfont{Latin Modern Math}
\usepackage{xurl,hyperref}
\hypersetup{colorlinks=true,urlcolor=blue,linkcolor=blue}
\newcommand{\sref}[2]{\hyperlink{source:#1:#2}{[S#2]}}
\newcommand{\cataloguescope}[1]{\paragraph{Mathematical question.}#1}
\begin{document}
% UnsolvedMath ID 8700056; source catalogue code AMR-086-0056
\setcounter{section}{8700055}
\section{Conjecture 4.14}\label{8700056}
{\small\sffamily UnsolvedMath ID 8700056 · Graph Theory}
\subsection{Definitions and mathematical statement}

\paragraph{Shared notation.}$\mathbb N=\{1,2,\ldots\}$, and $\mathbb Z,\mathbb Q,\mathbb R,\mathbb C$ denote the integers, rationals, reals and complex numbers. Any other conventions are specified locally.


For an algebraic number $\alpha$, choose its primitive irreducible polynomial $f_\alpha\in\mathbb Z[T]$ with positive leading coefficient $a$ and roots $\alpha^{(1)},\ldots,\alpha^{(d)}$. Set
\[
 M(\alpha)=a\prod_{j=1}^d\max\{1,|\alpha^{(j)}|\},\qquad
 \mathrm h(\alpha)=\frac{\log M(\alpha)}d,
\]
and let $H(\alpha)$ be the maximum absolute value of the coefficients of $f_\alpha$. These are respectively Mahler measure, absolute logarithmic height, and usual polynomial height.

Let $m\ge1$, let $\lambda_i\in\mathbb C$ satisfy $\alpha_i=e^{\lambda_i}\in\overline{\mathbb Q}^{\times}$, and let $\beta_0,\ldots,\beta_m\in\overline{\mathbb Q}$. Put $D=[\mathbb Q(\alpha_1,\ldots,\alpha_m,\beta_0,\ldots,\beta_m):\mathbb Q]$. Thus logarithms are chosen complex logarithms, with no restriction to a principal branch.
\paragraph{Statement recorded in the supplied source.}
Do there exist absolute constants $c_1,c_2>0$ such that, for every such tuple and every real $h$ satisfying
\[
h\ge \frac1D,\quad h\ge\max_i\mathrm h(\alpha_i),\quad h\ge\frac1D\max_i|\lambda_i|,\quad h\ge\max_{0\le j\le m}\mathrm h(\beta_j),
\]
both assertions hold?
\begin{enumerate}
\item If $\Lambda=\beta_0+\sum_{i=1}^m\beta_i\lambda_i\ne0$, then $|\Lambda|\ge e^{-c_1mD^2h}$.
\item If $\lambda_1,\ldots,\lambda_m$ are linearly independent over $\mathbb Q$, then $\sum_{i=1}^m|\lambda_i-\beta_i|\ge e^{-c_2mD^{1+1/m}h}$.
\end{enumerate} \sref{8700056}{2}

\subsection{Short English statement}
Are these two uniform quantitative lower bounds valid for nonzero algebraic linear forms in logarithms and for simultaneous approximation of independent logarithms by algebraic numbers?
\subsection{Sources}
\begin{description}
\item[\hypertarget{source:8700056:1}{S1}] ULAM.AI and the UnsolvedMath Contributors, UnsolvedMath: A Curated Collection of Open Mathematics Problems, record 8700056 (AMR-086-0056). CC BY 4.0. \url{https://huggingface.co/datasets/ulamai/UnsolvedMath}.
\item[\hypertarget{source:8700056:2}{S2}] Michel Waldschmidt, \emph{Open Diophantine Problems}, arXiv:math/0312440v2 (24 January 2004), Conjecture 4.14, pp. 40--41; height conventions, p. 35. \url{https://arxiv.org/abs/math/0312440v2}.
\end{description}
\label{end:8700056}
\end{document}
